\documentclass[10pt,a4paper]{amsart}
\usepackage[margin=19mm]{geometry}
\usepackage{amsmath,amssymb,mathtools}
\usepackage[scr=rsfs,cal=euler]{mathalfa}
\usepackage{xfrac}
\usepackage[unicode,hidelinks,bookmarks=false]{hyperref}
\newcommand{\ie}{i.e.\ }
\newcommand{\cf}{cf.\ }
\newcommand{\resp}{respectively\ }
\newcommand{\Subsection}{\subsection}
\providecommand{\nobreakdash}{\nobreak\hbox{-}}
\setlength{\emergencystretch}{2em}
\multlinegap0pt
\usepackage{amssymb}
\usepackage{xcolor}
\usepackage{enumerate}
\usepackage{dsfont}
\usepackage{bm}
\usepackage{tikz}
\usepackage{float}
\usepackage{csquotes}
\newtheorem{lem}{Lemma}
\newcommand{\R}{\mathbb{R}}
\title{Wave-front tracking for a quasi-linear scalar conservation law with hysteresis II: the case of Preisach}
\author{Fabio Bagagiolo\ \ and Stefan Moreti}
\date{Source: arXiv:2506.10087v2 (CC BY 4.0)}
\begin{document}
\maketitle

\noindent\emph{Source and licence.} Wave-front tracking for a quasi-linear scalar conservation law with hysteresis II: the case of Preisach. Version arXiv:2506.10087v2 (CC BY 4.0), distributed by arXiv under the Creative Commons Attribution 4.0 International licence, http://creativecommons.org/licenses/by/4.0/, as recorded in the arXiv licence metadata for this exact version and shown on the abstract page https://arxiv.org/abs/2506.10087v2. Full licence notice: \texttt{licenses/arxiv-2506.10087-CC-BY-4.0.txt}.\par\smallskip

\noindent\textit{Editorial scope.} Counted unit: the article's lem propPreis1 (source0.tex lines 508--510). The article's complete original proof of it is delivered with it (source0.tex lines 511--522, one \texttt{proof} environment of 12 lines). No mathematics is written, repaired or re-derived: the statement and the proof are the article's own lines, in the article's own notation, and no retained line is substituted or re-flowed.  No further source-local context is needed: every cross-reference of every retained line resolves inside the two delivered spans.  Nothing was omitted from any retained span, so the package carries no omission record.  No retained line contains a citation, so the wrapper carries no bibliography and no bibliography entry.  No retained line contains a figure environment, an image-inclusion command or drawing code, so no image is delivered and the package declares no asset dependency.  Editorial formatting only, in addition: an AMS \texttt{amsart} preamble that restates the article's own declarations and macros used by the retained lines, namely the environments \texttt{lem} on separate counters, together with the standard packages those lines need.  The retained environments print in this document as Lemma 1 in order of appearance, whereas the article prints the same items with the numbering of its own full document.  The complete native source of the article as distributed in the arXiv e-print for this exact version is bundled unchanged as \texttt{sources/179708/source0.tex}.
\begin{lem}\label{lemma: propPreis1}
    Consider a piecewise constant function $u:[0,T]\to\R$, consisting of a finite number of pieces, with its total variation $Var(u)\leq C$. Then there exists a function $u_{RMS}:[0,T]\to \R$ that is $L$-Lipschitz with $L$ depending on $C$ and $T$ such that $u_{RMS} (0)=u(0)$, $RMS(u;[0,T])=RMS(u_{RMS};[0,T])$ and $Var(u_{RMS})\leq C$.
\end{lem}

\begin{proof}
    Consider the reduced memory sequence of $u$, which will be of the form $\{\eta^1,\eta_2,\eta^3,\dots\}$ or $\{\eta_1,\eta^2,\eta_3,\dots\}$, where $\eta^i$ are maxima and $\eta_i$ are minima. We know that, denoting the elements of RMS by $v_i$ with $v_0=u(0)$, \[\sum_{i=1} |v_{i}-v_{i-1}|\leq Var(u)\leq C.\]
    Now define the sequence \[z_j:= \sum_{i=1}^{j}|v_i-v_{i-1}|\] so that $\{z_j\}_j$ is a monotonically increasing sequence bounded by $\tilde{C}\leq C$, with $\tilde{C} = \max z_j$ and set $z_0=u(0).$ By rescaling the points to the interval $[0,T]$, that is by setting
    \[
    x_i:= \frac{T}{\Tilde{C}} z_i,
    \] 
    we can define the function $u_{RMS}$ as the piecewise linear function that interpolates the points $(x_i,v_i)$. \par 
    It is clear that $RMS(u;[0,T])=RMS(u_{RMS};[0,T])$ and $Var(u_{RMS})\leq C$. \par
    Moreover, considering the modulus $m_i$ of the slopes of the segments corresponding to the linear parts we have that
    \[m_i = \frac{|v_i-v_{i+1}|}{|x_{i+1}-x_i|} = \frac{\tilde{C}}{T} \frac{|v_i-v_{i+1}|}{|z_{i+1}-z_i|} = \frac{\tilde{C}}{T} \leq \frac{C}{T},\]
    hence $u_{RMS}$ is $L-$lipshitz with $L\leq C/T$.
\end{proof}

\end{document}
